%!TEX root = ../../main.tex
\section{향후 과제}
\label{sec:disc-future}

다음 항목은 설계가 기록되어 있고 실행이 남은 것이다(부록 \ref{app:deferred}).

\begin{itemize}[leftmargin=1.5em]
  \item \textbf{동일 입력의 학습기 비교.} 같은 352열 입력, 전처리, 역할, 가중치에서 로지스틱과 LightGBM을 비교하여 학습 알고리즘의 선택이 예측에 미치는 독립적 영향을 잰다. 이 비교는 제한된 사양 사이의 것이며 모든 아키텍처의 우열로 일반화되지 않는다.
  \item \textbf{정의 변경의 점수 민감도와 재학습.} 일변량으로 바뀐 정의에서 동결 $q$·PT를 다시 점수화하고(frozen-score), $q$와 PT를 둘 다 다시 학습하는 모드(refit)를 같은 표본에서 수행한다.
  \item \textbf{사례 추적.} 질의 시각에서 사건 접두, $\Vhat$ 점수, $\SVI$, $q$ 예측, 후속 결과로 이어지는 사례별 추적을 사전에 정한 셀의 hash 표본에서 본문화한다.
  \item \textbf{미노출 표본의 점수 전용 평가.} 잠근 규칙(api\_patch $\ge$ 16.16, KR, 500 경기)대로 수집한 뒤 동결 $\Vhat$·$q$·PT를 점수 전용으로 적용한다.
  \item \textbf{외부 확률 어댑터.} 새 환경의 작은 보류 조각에서 적합한 저차원 단조 어댑터가 $\PTflex$에 적용한 같은 어댑터보다 Brier를 줄이는지, 가중치를 동결한 채 검정한다. 경쟁 가설(확률 척도 이동, 평가기 자체의 이동, 구성 이동)과 구별 예측을 함께 둔다.
  \item \textbf{참여 인원 셀.} $\nmin = 2 / 3 / 4 / 5+$ 셀 안에서 기존 예측표를 재점수화하여 코호트 안 이득의 분포를 기술한다.
  \item \textbf{공유 $\PTflexS$ 객체.} S arm 사이에 하나의 동결 $\PTflexS$를 재사용해도 대비가 변하지 않음을 확인하는 위생 항목이다.
\end{itemize}

다른 게임과의 통합, 모든 그래프·시계열 구조의 재실험, 실시간 교전 발생 탐지, 인과적 최적 행동 추천, 새 수작업 주석 연구, 연속형 $\dV$로의 주 목표 교체는 이 학위논문의 범위 밖이다.
